\section{Support, physical projection and parameter proofs}
\subsection{The classical support criterion}\label{app:support}
\begin{proof}[Proof of Proposition~\ref{prop:support}]
For any hull point $(e,c)$, linear minimization gives
$c+p^\top e\ge V(p)$. If $p\in\partial F(e(x))$, then
$F(e)\ge F(e(x))+p^\top(e-e(x))$. Adding and minimizing over the hull
gives
\[
 \CH\ge V(p)+F(e(x))-p^\top e(x)
       =c(x)+F(e(x))-r(x;p).
\]
This proves \eqref{eq:lower}. If $r(x;p)=0$, equality of physical and hull
values follows immediately.

Conversely let $D=\CH$. Compactness and lower semicontinuity give a
physical minimizer. The relative-interior qualification gives strong
Fenchel duality and an attained maximum
$\CH=V(p^*)-F^*(p^*)$. Under the alternative finite-minorant condition,
apply the same duality theorem to the finite minorant. It agrees with $F$
on hull loads and its conjugate is at least $F^*$. The corresponding
price gives a dual bound for $F$ at least $\CH$; weak duality forces
equality, hence attainment for $F$ as well.

For any physical minimizer $x$ and any attained hull-dual maximizer $p^*$,
\[
 \underbrace{c(x)+(p^*)^\top e(x)-V(p^*)}_{\ge0}
 +\underbrace{F(e(x))-(p^*)^\top e(x)+F^*(p^*)}_{\ge0}
 =D-\CH=0.
\]
Both terms vanish. Fenchel equality gives $p^*\in\partial F(e(x))$;
the first term gives support. For restricted economic prices the converse
requires an attained dual price in that subset, as stated.
Every dual maximizer satisfies the displayed unrestricted algebraic
conclusion; only the admissible maximizers give support under a restricted
economic definition.
\end{proof}

\subsection{Physical projection and the connector/fee region}\label{app:capacity}
\begin{proof}[Proofs of Theorems~\ref{thm:nominal} and~\ref{thm:capacity}]
With one bus, the state after A is 5 and the state after B is $x-10$.
Early charging has $x\le10$, so feasibility forces $x=10$. Terminal
replenishment is 20 and requires $K\ge20$. With two buses, A's bus has
terminal demand $15-x$, B's bus has demand 15, and $0\le x\le10$.
Serial charging at power $K$ fits the terminal hour exactly when
$30-x\le K$. These conditions also satisfy each battery's individual
headroom and inventory bounds. Thus for $K\ge20$ the two-bus early
interval is $[\ell,10]$, $\ell=30-K$; for $K<20$ neither partition
is feasible. At $K=20$ both projected loads are $(10,20)$ and their
costs differ only by $f$.

For $K>20$, strict convexity of $g$ gives the two-bus minimum at
$t=\max\{5,\ell\}$. The one-bus minimum is the fixed point 10,
proving \eqref{eq:capacityD}. Projecting complete plans to early energy
and operating cost gives the segment $\{(x,2f):\ell\le x\le10\}$
and the point $(10,f)$. Their convex hull has lowest cost
\[
 c_{\min}(x)=2f-\frac{f(x-\ell)}{10-\ell},\qquad\ell\le x\le10.
\]
The derivative of $c_{\min}(x)+g(x)$ vanishes at
$5+5f/(2d)$, $d=10-\ell$, proving \eqref{eq:capacityCH}.

For $f>0$, an interior hull minimizer has lower operating cost than every
two-bus physical plan at the same early load. Strict convexity excludes
the distinct one-bus point from attaining its objective. Therefore the
gap is positive exactly when $\ell<h<10$. Equivalently
$\tfrac25d(\ell-5)<f<2d$, together with $f>0$, which gives
\eqref{eq:capacityregion}. At an endpoint the physical endpoint attains
the hull minimum. For $f=0$, $h=t$, and the two-bus physical plan attains
the hull. Equality of the two physical branch minima gives exactly
$f=90-g(t)$; within each branch strict convexity fixes its minimizing
load. The affine-response formula \eqref{eq:fullresponse} follows by
minimizing over every point of those physical segments.

For $K=30,f=7$, the lower hull objective is
$104-27x/10+x^2/5=7591/80+(x-27/4)^2/5$. Physical branch values are 97
and 99. The one-bus gradient is $(6,4)$, and the complete two-bus
response at that price is the zero-early endpoint. Bills 147 and 134 give
regret 13. The mixture and operating-cost values in the text follow from
weight $27/40$ on the one-bus point.
\end{proof}

\subsection{A reserve/loss neighborhood}\label{app:reserve}
\begin{proof}[Proof of Proposition~\ref{prop:reserve}]
For reserve $0\le r\le5$, constant efficiency $0<\eta\le1$, early
power $P$, and terminal power $K$, define grid-energy quantities
\[
 T=30/\eta,\quad \ell=\max\{0,T-K\},\quad
 u=\min\{P,15/\eta\},\quad
 h=\max\{\ell,(10+r)/\eta\}.
\]
The two-bus early interval is $[\ell,u]$; the one-bus interval is
$[h,u]$ if nonempty. The lower terminal-capacity condition is
$T-x\le K$. For the one-bus plan, its state after B is $\eta x-10$,
so reserve imposes $x\ge(10+r)/\eta$. Charging only A's bus early
realizes every interval point; terminal demands fit serially, and the
individual bounds follow from the inventory headroom 15 and the stated
reserve range.

Write $G_T(x)=F(x,T-x)=G_T(m)+k(x-m)^2$, where $k=1/5$ and
$m=T/2-10$. Suppose
\begin{equation}\label{eq:strictreserve}
 \ell<m<h<u,\qquad f>0,\qquad f<2k(h-\ell)(h-m).
\end{equation}
On $[\ell,h]$ the hull's lowest operating cost is
$2f-f(x-\ell)/(h-\ell)$, obtained from the two-bus endpoint $\ell$
and one-bus endpoint $h$; on $[h,u]$ it is $f$. Put $d=h-\ell$.
The hull minimizer $x^*=m+f/(2kd)$ lies in $(m,h)$. Physical and hull
values are
\begin{align*}
 D&=G_T(m)+\min\{f+k(h-m)^2,2f\},\\
 \CH&=G_T(m)+2f-\frac{f(m-\ell)}d-\frac{f^2}{4kd^2}.
\end{align*}
Subtracting gives
\[
 \Delta=\min\left\{k(h-x^*)^2,
 \frac{f(m-\ell)}d+\frac{f^2}{4kd^2}\right\}>0.
\]
At $(r,\eta,P,K,f)=(1,19/20,12,30,7)$,
\[
 (\ell,m,h,u,d)=(30/19,110/19,220/19,12,10).
\]
In particular $u-h=8/19>0$ and $2kd(h-m)=440/19>7$.
Rational substitution gives the stated values. All inequalities in
\eqref{eq:strictreserve} are strict and the parameter functions are
continuous locally. They therefore persist in an open neighborhood of
$(r,\eta,P)$ with $K=30$ fixed.
\end{proof}

\section{Expansion and equipment proofs}
\subsection{The old-plan and dual-certificate statements}\label{app:expansion}
\begin{proof}[Proof of Proposition~\ref{prop:expansion}]
Set inclusion gives $V_Z(p)=\min\{V_X(p),V_Y(p)\}\le V_X(p)$, so
old-plan regret cannot decrease. By Proposition~\ref{prop:support}, no old
plan can become supported. If failure is erased, $D_Z=\CH_Z\le\CH_X<D_X$;
thus a new plan attains physical optimum and $D_Y\le\CH_X$. If instead
$\min\{D_X,D_Y\}>\CH_X$, then $D_Z>\CH_X\ge\CH_Z$, proving survival.

Under \eqref{eq:dualsafe}, $V_Z(p^*)=V_X(p^*)$, so
\[
 \CH_Z\ge V_Z(p^*)-F^*(p^*)=\CH_X\ge\CH_Z.
\]
For any new physical plan,
\begin{align*}
 c(y)+F(e(y))-\CH_X
 ={}&[c(y)+(p^*)^\top e(y)-V_X(p^*)]\\
 &+[F(e(y))-(p^*)^\top e(y)+F^*(p^*)].
\end{align*}
Both brackets are nonnegative. Equality holds exactly under the fleet and
supply equalities stated in the proposition. Compactness ensures that if
$D_Z=\CH_X$, some physical plan attains that equality; it cannot be old.
With a uniform fleet margin $\varepsilon$, every new physical plan lies
at least $\varepsilon$ above $\CH_X$, while every old one lies at least
$\Delta_X$ above it. This gives the quantitative bound. Strict convexity
of $F$ makes the minimizer of $F(e)-(p^*)^\top e$ unique on its convex
domain whenever attained, proving the projected-point interpretation.
\end{proof}

\begin{proof}[Proof of Theorem~\ref{thm:movement}]
Each added plan buys total energy $30+w$, $w\ge0$, because its initial
and final battery inventories agree. Its early purchase $u$ is at most 10.
For two used buses, the old $(u,30-u)$ plan of cost 14 is feasible and
componentwise no larger in load or cost. For one bus, completion of both
services by hour 3 requires $20+u\ge30$ even before any additional
nonnegative movement consumption. Thus $u=10$, and the old one-bus
plan dominates it. No added feasible plan can reduce the physical
objective because $F$ is coordinatewise nondecreasing.

Every hull point is a finite convex combination of physical points.
Replace each added component by its old dominator; its mixture has no
larger cost or load in any coordinate. Consequently it has no larger
convex supply objective. This gives $\CH_X\le\CH_Z$, and inclusion
gives the reverse inequality. The same inclusion/dominance argument
gives $D_X=D_Z$. This proof does not require the movement timing to be
favorable: it concerns every added plan that remains feasible.
\end{proof}

\subsection{Battery enlargement and fee thresholds}\label{app:battery}
\begin{proof}[Proof of Theorem~\ref{thm:battery} and
\eqref{eq:B25D}--\eqref{eq:B25CH}]
The one-new-bus state sequence imposes $x\ge30-B=a$ and $x\le10$.
Every point of $[a,10]$ is feasible and fully replenished. With two
buses, the A bus's headroom after its service is 15 for either type;
B's bus starts full before service and cannot buy early energy. Serial
terminal charging gives every $x\in[0,10]$ for every type assignment.
Since $q\ge7$, all mixed/new two-bus projected points are dominated by
the old two-bus interval. Therefore
\[
 D_Z=\min\{97,q+g(t)\},\qquad t=\max\{5,a\}.
\]
Every supported physical plan is a physical planner optimum. No projected
old plan can be rescued, and a dominated two-bus addition cannot have
zero regret: at equal cost it has an unsupported old counterpart; at
larger cost it has a strictly cheaper old counterpart at every price.
It suffices to test the new branch's unique physical minimizer $t$.

If $a\le5$, $t=5$, so $p=(5,5)$. All plans buy 30 units; comparison
with the old one-bus plan is simply $q\le7$, hence $q=7$.
If $a>5$, $t=a$ and $p_1-p_2=2a/5-2>0$. Within the new branch,
the endpoint $a$ minimizes the fixed-price objective. The best old
two-bus plan is its zero-early endpoint. Comparison against it gives
\[
 q\le14-a(p_1-p_2)=14+2a-\tfrac25a^2=A(a).
\]
Comparison against the old one-bus point gives
\[
 q\le7+(10-a)(p_1-p_2)=C(a).
\]
These comparisons include all physical alternatives by affine minimization
on every branch. They prove \eqref{eq:batteryregion}. For $a>5$,
$C(a)\ge7$ and $A(a)\ge7$ exactly when
$a\le(5+\sqrt{95})/2$ in the relevant domain; also $A(a)-C(a)=27-4a$.
At $q=7$, including $a\le5$, this proves \eqref{eq:batterythreshold}.

At fee 7 below the battery threshold, the hull's lower boundary joins
$(0,14)$ to $(a,7)$ and is then flat to 10. The hull objective on the
first edge is $14-7x/a+g(x)$ and is minimized at
$t_h=5+35/(2a)<a$. The physical optimum has early energy $a$;
its gap is $(a-t_h)^2/5$. At and above the threshold the physical
point $(\max\{5,a\},7)$ attains the hull minimum. Thus the boundary is
derived on the whole continuous battery domain, not inferred from a scan.

For $B=25$, the added endpoint is $(5,q)$. It is below the old lower
hull $14-0.7x$ exactly when $q<21/2$. In that regime the new lower
boundary joins $(0,14)$ to $(5,q)$ to $(10,7)$. The objective decreases
along the first edge. On the second it is minimized at
$x=5+(q-7)/2$ with operating cost $q-(q-7)^2/10$. Its value is
$85+q-(q-7)^2/20$. At $q=21/2$ this equals $7591/80$; above it the
whole new segment is dominated by the old hull. This proves the formulas.
For $q=8$, substitution gives $\CH_Z=1859/20$ and $D_Z=93$.
For $q=12$, both physical minima are 97. At the new load's price $(5,5)$,
the old one-bus bill is cheaper by 5. At the old load's price $(6,4)$,
the old zero-early two-bus bill remains 134; every new one-bus bill is
at least $q+130\ge137$. Thus the old plan's regret is 13, including
at the fee 12 tie. All quantities use fixed posted prices.
\end{proof}

\section{Replication proof}\label{app:replication}
\begin{proof}[Proof of Theorem~\ref{thm:replication} and
Corollary~\ref{cor:individual}]
For $m$ paired buses, the complete projected physical set has early energy
$E\in[10m,10n]$, late energy $30n-E$ and operating cost $14n-7m$.
The connector assignment in the text realizes every such early total:
paired buses take 10 each, and unpaired A buses distribute the remainder.
Convexification has lower operating boundary $14n-0.7E$. Substitution
in \eqref{eq:scaledF} gives the hull minimizer $E=27n/4$ and value
$7591n/80$.

For fixed integer $m\ge n/2$, supply is minimized at $E=10m$ and total
cost is
\[
 104n-27m+\frac{20m^2}{n}
 =\frac{7591n}{80}+\frac{20}{n}(m-27n/40)^2.
\]
For $m\le n/2$, the best early energy is $5n$ and branch cost is
$99n-7m$, decreasing with $m$. For odd $n$, moving from
$\lfloor n/2\rfloor$ to $\lceil n/2\rceil$ changes the minimized
cost by $-7+5/n<0$. For even $n$ the branch formulas meet at $n/2$.
Thus an integer nearest to $27n/40$ is optimal, including both at a tie,
and $|\delta|\le1/2$. This proves \eqref{eq:scaledgap}.

At that optimum the own-load marginal price is
\[
 p=(4+2m/n,\ 6-2m/n).
\]
The optimal $m$ has $m/n\ge1/2$, so a cheapest two-bus pair response
buys zero early (the whole interval ties when $m/n=1/2$).
The one-bus private objective minus that two-bus
objective is $40\delta/n$. If $\delta\ge0$, precisely the $m$
paired choices are disadvantaged; otherwise the $n-m$ unpaired choices
are. This proves \eqref{eq:scaledregret}. With one owner, all choices can
be replaced simultaneously. To compare separate ownership, first implement
the optimum using $m$ fixed-pair owners with one bus and $n-m$ with two.
Each one-bus owner buys 10 early; each two-bus owner buys zero early at the
planner optimum, giving the same total load and operating cost as the
aggregate optimum even if its literal matching crossed the new ownership
labels. Reserved connectors make each unilateral replacement feasible.
The identical projected opportunity sets give the same complete-fleet
response value, while each owner's gain is at most $20/n$. Summing the
disadvantaged owners' gains yields the whole-owner expression.

For $n=40k+1$, the optimum is $m=27k+1$, so $\delta=13/40$;
substitution yields the displayed subsequence. If $40$ divides $n$,
$\delta=0$. At $n=20$, the tied $m=13,14$ have $\delta=-1/2,1/2$,
giving regrets 7 and 14 respectively. All ties concern physical plans;
no fractional bus schedule is being implemented.
\end{proof}
